\chapter{Endoscopic Sleeve Gastroplasty}

\section*{Overview}
Endoscopic sleeve gastroplasty (ESG) is a minimally invasive procedure in which the stomach is reduced in size with internal sutures placed through an endoscope, restricting food intake and promoting weight loss in patients with obesity.
\section*{Alternative Names}
ESG; endoscopic suturing of the stomach; sleeve gastroplasty without surgery; gastric sleeve endoscopy
\section*{Principle}
An endoscopic suturing device places a series of full-thickness stitches in the gastric wall, creating a narrow tubular channel along the lesser curvature and markedly reducing the stomach's capacity, which promotes early satiety and reduced intake.
\section*{History}
Endoscopic sleeve gastroplasty was developed in the 2010s as a less invasive alternative to laparoscopic sleeve gastrectomy and has gained acceptance as a primary or adjunctive weight-loss procedure.
\section*{Why It Is Done}
Ordered for patients with class I or II obesity or with obesity-related complications who have failed lifestyle measures, and who prefer or qualify for an endoscopic rather than a surgical weight-loss procedure.
\section*{Is This a Primary or Secondary Test or Procedure}
A primary therapeutic procedure for weight management in selected patients.
\section*{How It Is Performed}
Under general anesthesia, an endoscope with a suturing device is passed into the stomach. Sutures are placed along the gastric wall to fashion the sleeve, after which the scope is removed. The procedure usually lasts under an hour.
\section*{Preparation}
Preoperative evaluation including nutritional assessment and exclusion of gastric pathology, review of medications, and a clear liquid diet before the procedure.
\section*{Contraindications and Precautions}
Contraindications include severe gastroparesis, gastric ulceration, prior gastric surgery, uncontrolled psychiatric disease, and bleeding disorders. Precautions include pregnancy and large hiatus hernias.
\section*{Risks}
Nausea and vomiting after the procedure, pain, and rarely bleeding, perforation, suture disruption, or gastroesophageal reflux. Weight regain may occur over time.
\section*{Aftercare and Recovery}
The patient follows a staged diet, progressing from liquids to soft foods over several weeks, and is monitored with regular follow-up. Nausea is managed with antiemetics.
\section*{Outcome and Follow-up}
Weight loss typically begins in the first weeks and continues over months. The degree of loss is assessed against the patient's goals and the expected range for the procedure.
\section*{Expected Results}
Significant and sustained weight loss with reduction of obesity-related conditions and no serious complications.
\section*{Follow-up}
Ongoing nutritional monitoring and lifestyle support are required, and repeated or additional procedures may be considered if weight loss is insufficient or weight is regained.
\section*{When to Seek Medical Care}
Follow the procedure team's specific instructions. Seek urgent medical assessment for severe or worsening pain, uncontrolled bleeding, fever or signs of infection, trouble breathing, chest pain, fainting, new weakness or confusion, inability to urinate, or any rapidly worsening symptom. The appropriate warning signs depend on the procedure and the patient's medical condition.

\section*{References}
\begin{enumerate}
  \item MedlinePlus. Weight-loss procedures. U.S. National Library of Medicine; 2025.
  \item Current Medical Diagnosis and Treatment. McGraw-Hill; 2025.
\end{enumerate}

\clearpage
